\section*{ВВЕДЕНИЕ}
\addcontentsline{toc}{section}{ВВЕДЕНИЕ}

Организация записи клиентов является важной частью работы частных специалистов, оказывающих услуги по предварительной записи. К таким специалистам относятся мастера индустрии красоты, преподаватели, консультанты, тренеры и другие исполнители, работа которых связана с фиксированными временными интервалами.

На практике запись часто ведется вручную: через телефонные звонки, сообщения в мессенджерах, бумажный журнал или электронные таблицы. Такой подход требует постоянного участия мастера, усложняет контроль свободного времени и повышает риск двойного бронирования. Дополнительную сложность создает необходимость учитывать личный календарь мастера, своевременно уведомлять клиента об изменениях и хранить историю взаимодействия в едином месте.

Для решения этих задач целесообразно использовать web-платформу онлайн-бронирования. Такая система позволяет клиенту самостоятельно выбрать доступный слот, а мастеру -- управлять расписанием, приглашениями и записями через личный кабинет. При наличии интеграции с внешним календарем система может исключать занятые интервалы из доступного расписания, а при наличии канала уведомлений -- автоматически информировать участников процесса о создании, изменении или отмене записи.

Разрабатываемая платформа Easyappoint предназначена для автоматизации процесса записи клиентов к частному мастеру. Система объединяет web-интерфейс, серверную часть, базу данных, интеграцию с Google Calendar и сервис Telegram-уведомлений. Событийный обмен между компонентами используется для доставки уведомлений и повышения связности бизнес-процесса без жесткой зависимости между сервисами.

\emph{Цель настоящей работы} -- разработка платформы онлайн-бронирования услуг на базе событийно-ориентированной микросервисной архитектуры. Для достижения поставленной цели необходимо решить \emph{следующие задачи:}
\begin{itemize}
\item провести анализ предметной области онлайн-бронирования услуг;
\item определить требования к функциональности, интерфейсу, надежности и безопасности платформы;
\item спроектировать архитектуру программной системы, модель данных и интеграционные взаимодействия;
\item реализовать серверную часть, клиентское web-приложение и сервис уведомлений;
\item выполнить тестирование ключевых сценариев записи, отмены, изменения бронирований и отправки уведомлений.
\end{itemize}

\emph{Структура и объем работы.} Отчет состоит из введения, 4 разделов основной части, заключения, списка использованных источников, 2 приложений. Текст выпускной квалификационной работы равен \formbytotal{lastpage}{страниц}{е}{ам}{ам}.

\emph{Во введении} сформулирована цель работы, поставлены задачи разработки, описана структура работы, приведено краткое содержание каждого из разделов.

\emph{В первом разделе} приводится анализ предметной области онлайн-бронирования услуг, рассматриваются проблемы ручной записи и обосновывается необходимость автоматизации.

\emph{Во втором разделе} на стадии технического задания приводятся требования к разрабатываемой платформе.

\emph{В третьем разделе} на стадии технического проектирования представлены проектные решения для архитектуры, базы данных и основных компонентов платформы.

\emph{В четвертом разделе} приводится описание программной реализации, компонентов и тестирования разработанной платформы.

В заключении излагаются основные результаты работы, полученные в ходе разработки.

В приложении А представлен графический материал. В приложении Б представлены фрагменты исходного кода.
